\scanpage{93}
\mspage{92}

développements actuels ne dépendent de $\unsure{I|F}$ que par l'intermédiaire des Strates $\unsure{M|\mathcal M}$, $\leq$, $\ll$, $\unsure{\ngeq|\gtrless}$.

\deleted{Il sera commode de \unsure{désidentifier}, et d'introduire un ens. d'indices ordonné $I$, pour \unsure{indexer} $\widetilde F$}

Pour toute partie finie $A$ de $\widetilde F=\Phi$, posons
\[
\begin{aligned}
S_A&=\operatorname{Supp}_{M}\left(\bigcup_{X\in A}X^\circ\right)\quad\in\overline\Sigma_M\\
&=\underset{X\in A}{\operatorname{Sup}_{\Sigma_M}}\operatorname{Supp}(X^\circ).
\end{aligned}
\]
\margin{$\deleted{\unsure{\Rightarrow}}$}

(\uline{NB} Si on admet At supp. 1, on a $\operatorname{Supp}X^\circ=X^\circ$, i.e. $X^\circ$ est déjà une partie-support de $M$.)

Par définition et par associativité des Sup, on aura donc, si $A=\bigcup_i A_i$
\[
S_A=\operatorname{Sup}_{\deleted{\unsure{i}}\Sigma_M}S_{A_i}.
\]
On a ainsi une application
\diagram{\begin{tikzcd}[ampersand replacement=\&,column sep=large]
\unsure{\mathscr P|\mathcal P}(\widetilde F\unsure{]|)}\arrow[r]\&\Sigma_M\ ;\&A\arrow[r,mapsto]\&S_A
\end{tikzcd}}
de l'ens. ordonné des parties de $\widetilde F$ dans l'ens \deleted{des} ordonné des supports pour $M$, application évidemment croissante.
